% ===== Block H: PDF p.35–37 | 原文 §3.3 及其之前的 p.35 段落 | 唯一编号公式 3.49–3.60 =====
% [原文 p.35]
% [NOTE: 按 Lead 裁决，本块从 PDF p.35 页面顶部开始。首句 “where $\delta_m=...$” 是 p34 末 (3.48) 的续句，原文此处为斜体（承接上一块的定理/命题体），故用 \emph 排出。原文 §3.3 标题位于 p35 中部，按 Lead 指示由本文件就地写入 \subsection。]

\emph{其中 $\delta_m=\frac12 j_*\left(\sum_{k+\ell=m-1}\psi^k(\psi')^\ell\right)$，$j$ 是全部余维一边界层（即恰有一条边的稳定图）的包含映射。类 $\psi$ 与 $\psi'$ 是结点处的两个 $\psi$ 类。}

Givental 作用之所以极其有力，有两个原因。第一，我们很快就会看到，它给出了一种递归计算 CohFT 关联函数的方法。第二，它或许还能产生上同调中的关系！以 Mumford 公式为例。由几何上的理由我们知道，Hodge 类 $\Lambda(t)$ 在次数 $d>g$ 时消失（它是秩为 $g$ 的丛的 Chern 多项式）。另一方面，$\Lambda(t)$ 的 Mumford 公式在任意次数都给出某个类。把 Mumford 公式在复次数 $d$ 中的分量（即式 (3.48) 右端中 $t^d$ 的系数）记为 $\mathcal H^d_{g,n}$，我们便得到如下的重言关系：对每个 $d>g$，在 $H^{2d}(\Mc_{g,n},\Q)$ 中都有 $\mathcal H^d_{g,n}=0$。这类重言关系的第一个非平凡例子是亏格 $0$ 中次数为 $1$ 的关系：

% [Tutor Note: 这里的论证是“两面夹击”：一方面 $\Lambda(t)$ 是秩 $g$ 丛的 Chern 多项式，故次数 $>g$ 的分量为零；另一方面 Mumford 公式（式 (3.48)）把每个次数分量 $\mathcal H^d_{g,n}$ 写成显式的重言类。两者对照即得 $\mathcal H^d_{g,n}=0$；注意 $\mathcal H^d_{g,n}$ 落在 $H^{2d}(\Mc_{g,n},\Q)$，而不是整个上同调。]

\begin{equation}
  \mathcal H^1_{0,n}=\kappa_1-\sum_{i=1}^{n}\psi_i+\delta_1=0 \qquad \text{in } H^2(\Mc_{0,n},\Q).
  \tag{3.49}
\end{equation}

% [Tutor Note: (3.49) 中的 $\delta_1$ 是 p.35 开头给出的 $\delta_m=\frac12 j_*(\sum_{k+\ell=m-1}\psi^k(\psi')^\ell)$ 在 $m=1$ 时的情形（$j$ 为所有余维一边界层之并），它不是某些文献中特指某一分离型边界除子的 $\delta_1$，两套记号不可混用。核对：$n=3$ 时 $\Mc_{0,3}$ 为单点、$H^2(\Mc_{0,3},\Q)=0$，关系平凡成立；第一个非平凡情形是 $n=4$。]

Pixton–Pandharipande–Zvonkine [PPZ15] 在 Witten 3-spin 类的情形中运用了这一论证，证明了上同调中所有已知的关系。

\begin{exercise}[3.3]
利用 Mumford 公式证明 $\Lambda(t)\Lambda(-t)=1$。这有时被称为 Mumford 关系。并由此推出关系 $\lambda_g^2=0$。
\end{exercise}

% [Tutor Note: 由 $\Lambda(t)\Lambda(-t)=1$ 推 $\lambda_g^2=0$，只需看 $t^{2g}$ 的系数：$\lambda_i=0$（$i>g$），故该系数中只剩 $i=g$ 的一项 $(-1)^g\lambda_g^2$，于是 $\lambda_g^2=0$。特别地 $g=1$ 时 $\Lambda(t)=1+\lambda_1t$，直接得 $\lambda_1^2=0$。]

Givental 作用极其有价值的另一个原因是它的适用范围，这一结果由 Teleman [Tel12] 证明。Teleman 证明了：凡是底层的量子积为半单的 CohFT，都包含在平凡 CohFT 于 Givental 作用下的轨道中。在附加的齐次性条件下，他给出了显式计算旋转矩阵与平移矩阵的算法。

\begin{theorem}[3.8（Teleman 分类）]
设 $\Omega$ 是 $(V,\eta)$ 上的一个 CohFT。若 $\Omega$ 半单且齐次，则存在显式的 $R$ 与 $T$，使得
\begin{equation}
  \Omega_{g,n}=RTw_{g,n},
  \tag{3.50}
\end{equation}
其中 $w_{g,n}=\Omega_{g,n}|_{H^0(\Mc_{g,n},\Q)}$ 是相应的二维拓扑场论（2D TFT）。若 $\Omega$ 半单（但不齐次），则存在 $R$ 与 $T$ 使上式成立，但它们被确定到相差一个对角歧义（diagonal ambiguity）为止。
\end{theorem}

% [Tutor Note: $w_{g,n}=\Omega_{g,n}|_{H^0(\Mc_{g,n},\Q)}$ 取的是 CohFT 的次数 $0$ 分量；由 (3.52) 可见，对平凡的副本这一分量恰为纯数 $\delta_{\mu_1,\dots,\mu_n}/(t^{\mu_i})^{2g-2+n}$，即“2D TFT 的关联函数”是一族这样的数。定理 3.8 中的 $R,T$ 在 §3.3 中被具体化：$R$ 见 (3.53)，$T$ 见 (3.54)。]

换言之，Teleman 定理把所有半单且齐次的 CohFT 刻画为半单 2D TFT 在 Givental 作用下的轨道。相当漂亮！

% [NOTE: 原文此处为 run-in 式标题：“3.3. Connection to topological recursion. Givental's action provides a recursive construction of CohFTs.”（标题与正文同段）。按 Lead 指示就地写 \subsection；紧接标题后的这句话作为小节首段。]

\subsection{与拓扑递归的联系（Connection to topological recursion）}

Givental 作用给出了一种 CohFT 的递归构造。

由于平凡 CohFT 的关联函数可由拓扑递归递推地算出，一个自然的问题随之而来：能否递归地计算平凡 CohFT 的 Givental 轨道中得到的全部关联函数？答案是肯定的，而且这与拓扑递归的理论有着美妙的联系。

% [Tutor Note: 平凡 CohFT 的关联函数就是 $\psi$ 类相交数 $\int_{\Mc_{g,n}}\prod_i\psi_i^{d_i}$，而 Witten 猜想／Kontsevich 定理（见 §2.3）说它们可由拓扑递归计算。于是“Givental 轨道中的关联函数能否递归计算”这一问题，正对应下面的定理 3.9：拓扑递归的 $\omega_{g,n}$ 给出轨道中 CohFT 的关联函数。]

% [原文 p.36]

考虑一条带有 $r$ 个单分歧点（simple ramification points）的谱曲线 $\mathcal S=(\Sigma,x,y,B)$。围绕分歧点 $\mu$ 选取局部坐标 $\zeta_\mu$，使得 $x=\zeta_\mu^2+x(\mu)$。考虑辅助函数 $\xi^\mu$ 以及与之相伴的亚纯微分 $\dd\xi^{\mu,k}$，它们定义为

\begin{equation}
  \xi^\mu(z)=\int^{z}\frac{B(w,\cdot)}{\dd\zeta_\mu(w)}\Big|_{w=\mu},
  \qquad
  \dd\xi^{\mu,k}(z)=\dd\left(\left(-\frac{1}{\zeta_\mu}\frac{\dd}{\dd\zeta_\mu}\right)^{k}\xi^\mu(z)\right).
  \tag{3.51}
\end{equation}

% [Tutor Note: (3.51) 第一式的积分 $\int^z$ 没有写下限，故 $\xi^\mu$ 只定义到相差一个常数；但后文只用到 $\dd\xi^{\mu,k}$（(3.53) 中还有 $\dd\xi^\mu$），常数不影响结果。]

% [Tutor Note: 核对 (3.51) 在分歧点处的极点阶：局部取 $B(w,z)=\frac{\dd w\,\dd z}{(w-z)^2}$，在 $w=\mu$ 处得 $B(\mu,z)/\dd\zeta_\mu(\mu)=\dd z/\zeta_\mu(z)^2$，故 $\xi^\mu(z)\sim-1/\zeta_\mu(z)$。于是 $(-\zeta_\mu^{-1}\dd/\dd\zeta_\mu)^k\xi^\mu\sim\zeta_\mu^{-1-2k}$，而 $\dd\xi^{\mu,k}$ 在 $\mu$ 处有 $2k+2$ 阶极点——这正是 (3.56) 中 $\tau_{\mu_i,d_i}$ 与 $\dd\xi^{\mu_i,d_i}$ 配对的基。]

令 $t^\mu=-2\frac{\dd y(z)}{\dd\zeta_\mu(z)}\Big|_{z=\mu}$。在 $V=\bigoplus_{\mu=1}^{r}\C\cdot e_\mu$ 上定义（平凡 CohFT 的 $r$ 个副本）CohFT\footnote{本节余下部分中，我们在 $\C$ 上而非 $\Q$ 上工作。}：令 $\eta(e_\mu,e_\nu)=\delta_{\mu,\nu}$，并规定

\begin{equation}
  w_{g;\mu_1,\dots,\mu_n}=\frac{\delta_{\mu_1,\dots,\mu_n}}{(t^{\mu_i})^{2g-2+n}}.
  \tag{3.52}
\end{equation}

% [Tutor Note: 原文此处的脚注（“在 $\C$ 上而非 $\Q$ 上工作”）在底本中是全文连续编号的第 4 条；译本用 LaTeX 自动编号，故显示的数字会与原文不同。这也解释了为什么 (3.49) 与定理 3.8 在 $\Q$ 上叙述、而以下 (3.55) 的靶空间改用 $\C$ 系数。]

% [Tutor Note: (3.52) 分母写作 $t^{\mu_i}$（上标 $\mu$、下标 $i$），严格说应为 $t^{\mu_1}$ 或 $t^\mu$；因分子 $\delta_{\mu_1,\dots,\mu_n}$ 已强制诸 $\mu_i$ 相等，表达式本身无歧义（已登记 ERRATA）。]

% [Tutor Note: (3.52) 的指数 $2g-2+n$ 可以这样核对：把 $w$ 代入 Givental 作用的稳定图求和（即把 $t^\mu$ 换成 $T^\mu(u)$），每个顶点 $v$ 贡献 $u^{2g_v-2+n_v}$；由 $\sum_v(2g_v-2+n_v)=2g-2+n$（用到 $\sum_v g_v=g-E+V-1$ 与 $\sum_v n_v=n+2E$），总幂次恰为 $u^{2g-2+n}$，与分母指数一致。]

如下定义旋转矩阵 $R$ 与平移 $T$：

\begin{equation}
  R^\nu_\mu(u)=-\sqrt{\frac{u}{2\pi}}\int_{\gamma_\nu}\ee^{-\frac{x-x(\nu)}{2u}}\dd\xi^\mu,
  \tag{3.53}
\end{equation}

\begin{equation}
  T^\mu(u)=\left(ut^\mu+\frac{1}{\sqrt{2\pi u}}\int_{\gamma_\mu}\ee^{-\frac{x-x(\mu)}{2u}}\omega_{0,1}\right).
  \tag{3.54}
\end{equation}

% [Tutor Note: (3.53) 与 (3.54) 中指数里的变量极易抄错：$R^\nu_\mu$ 里是 $x(\nu)$、路径为 $\gamma_\nu$；$T^\mu$ 里是 $x(\mu)$、路径为 $\gamma_\mu$。两处的积分都按 $u\to0$ 的渐近意义理解，见下段说明。]

% [Tutor Note: 为什么 $\gamma_\mu$ 可以取 $\zeta_\mu$-平面内的实轴：由 $x=\zeta_\mu^2+x(\mu)$ 得 $x-x(\mu)=\zeta_\mu^2$，沿实轴有 $\zeta_\mu^2\ge0$，故当 $|\zeta_\mu|\to\infty$ 时 $\ee^{-(x-x(\mu))/(2u)}$（$u>0$）指数衰减，这正是“最速下降”的初等情形。]

这里 $\gamma_\mu$ 是从分歧点 $\mu$ 出发的、关于 $x(z)$ 的形式最速下降路径；局部上它可以取在 $\zeta_\mu$-平面内的实轴上。此外，这些等式应理解为关于 $u$ 的形式幂级数之间的等式，其中右端取的是 $u\to0$ 时的渐近展开。

于是，通过 Givental 作用，我们可以由数据 $(w,R,T)$、按照 3.2 小节中所解释的对稳定图的求和，定义一个 CohFT

\begin{equation}
  \Omega_{g,n}=RTw_{g,n}\colon V^{\otimes n}\longrightarrow H^{2\bullet}(\Mc_{g,n},\C)
  \tag{3.55}
\end{equation}

它与拓扑递归关联函数之间的联系由下述定理给出 [Eyn14; DOSS14]。

% [Tutor Note: (3.55) 中的 $H^{2\bullet}$ 表示偶次上同调 $\bigoplus_{d}H^{2d}$；CohFT 的关联函数取值都在偶次部分，故靶空间写成 $H^{2\bullet}(\Mc_{g,n},\C)$ 而非整个上同调。另 (3.55) 只是把 (3.50) 中的 $R,T,w$ 按 (3.52)–(3.54) 具体化。]

\begin{theorem}[3.9（CohFT/TR 对应）]
固定一条紧谱曲线 $\mathcal S=(\Sigma,x,y,\omega_{0,2})$，并按式 (3.55) 定义 CohFT $\Omega$。则拓扑递归的关联函数计算 CohFT 的关联函数：
\begin{equation}
  \omega_{g,n}(z_1,\dots,z_n)=\left\langle\tau_{\mu_1,d_1}\cdots\tau_{\mu_n,d_n}\right\rangle_g^{\Omega}\dd\xi^{\mu_1,d_1}(z_1)\cdots\dd\xi^{\mu_n,d_n}(z_n).
  \tag{3.56}
\end{equation}
反之，若给定一个位于半单 2D TFT 的 Givental 轨道中的 CohFT，则可通过式 (3.53) 与 (3.54) 定义一条（局部的）谱曲线，它按式 (3.56) 那样计算关联函数。
\end{theorem}

% [Tutor Note: (3.56) 中 $\langle\tau_{\mu_1,d_1}\cdots\tau_{\mu_n,d_n}\rangle_g^{\Omega}$ 是 CohFT $\Omega$ 的关联函数（对 §3.1 中 $\langle\tau_{d_1}\cdots\rangle_{g,n}$ 附加上协变指标 $\mu_i$）；等式两边都是 $z_1,\dots,z_n$ 的微分形式，$\dd\xi^{\mu_i,d_i}$ 的极点阶 $2d_i+2$ 与拓扑递归 $\omega_{g,n}$ 的极点结构相符。]

% [Tutor Note: 定理 3.9 的陈述里谱曲线数据从 $(\Sigma,x,y,B)$ 换成了 $(\Sigma,x,y,\omega_{0,2})$：$B$ 就是 Bergman 核，即 $(g,n)=(0,2)$ 的关联函数 $\omega_{0,2}$，它是拓扑递归的初值数据（另一项是 $\omega_{0,1}$），也正是表 2 中“edge contribution”一行。]

简而言之，CohFT 与拓扑递归之间的对应可以概括为表 2。

\begin{exercise}[3.4]
证明：与下述谱曲线
\begin{equation}
  \left(\PP^1,\ x(z)=\frac{z^2}{2},\ y(z)=\frac{\sin(2\pi z)}{2\pi},\ \omega_{0,2}(z_1,z_2)=\frac{\dd z_1\dd z_2}{(z_1-z_2)^2}\right)
  \tag{3.57}
\end{equation}
相伴的 CohFT 是 Weil–Petersson CohFT $\exp(2\pi^2\kappa_1)$。另见习题 3.2；证明可参看 [Eyn, section 8.1]。
\end{exercise}

% [Tutor Note: 这是定理 3.9 反向陈述的具体例子：把 Weil–Petersson CohFT $\exp(2\pi^2\kappa_1)$ 对应的局部谱曲线写出来。核对分歧点：$x(z)=z^2/2$ 在 $z=0$ 有唯一的单分歧点，故 $r=1$，与 $\dim V=1$（表 2 第一行）相合。]

% [原文 p.37]
% [TODO-LEAD: 下面 Table 2 按 STYLE_GUIDE §7.2 用了 \caption*{...}，但当前 src/preamble.tex 未加载 caption 宏包（ctexart 不自动提供 \caption*），编译前请在 preamble 加 \usepackage{caption}；表 1/表 3 同理。]

\begin{table}[htbp]\centering\small
\caption*{表 2：CohFT 与拓扑递归数据之间的对应。}
\begin{tabular}{c|c}
\toprule
CohFT & 拓扑递归\\
\midrule
$\dim(V)$ & 分歧点个数\\
平凡 CohFT & $\dfrac{\dd y}{\dd\zeta}$\\
平移 & $\omega_{0,1}$\\
旋转 & $\dd\xi$\\
边贡献 & $\omega_{0,2}$\\
\bottomrule
\end{tabular}
\end{table}

% [Tutor Note: 表 2 是两边数据的对照表：$\dim(V)$（CohFT 态空间的维数，即幂等元个数）对应谱曲线的分歧点数 $r$；“trivial CohFT” 对应 $\dd y/\dd\zeta$（$t^\mu=-2\dd y/\dd\zeta_\mu|_{z=\mu}$ 的来源）；translation 对应 $\omega_{0,1}$（即 (3.54) 中被积的初值数据）；rotation 对应 $\dd\xi$（即 (3.53) 中的积分核）；edge contribution 对应 $\omega_{0,2}$（Bergman 核，即 Givental 图求和中每条边的贡献）。表中未单列 $w_{g,n}$，它由 $\dd y/\dd\zeta$ 这一行经 (3.52) 决定。]

\begin{exercise}[3.5]
证明：与下述谱曲线
\begin{equation}
  \left(\PP^1,\ x(z)=-f\log(z)-\log(1-z),\ y(z)=-\log(z),\ \omega_{0,2}(z_1,z_2)=\frac{\dd z_1\dd z_2}{(z_1-z_2)^2}\right).
  \tag{3.58}
\end{equation}
相伴的 CohFT 是三重 Hodge 类 $\Lambda(1)\Lambda(f)\Lambda(-f-1)$。证明见 [Eyn, section 7]。这是（带框的）拓扑顶点（topological vertex）[MV02; LLZ03; OP04] 背后的 CohFT，而三重 Hodge 类的拓扑递归公式不是别的，正是顶点的 BKMP 重构猜想（BKMP remodelling conjecture）。大 framing 极限（large framing limit）可还原出 [BM08] 中所谓的 Lambert 曲线，后者计算 Hurwitz 数（参看 [Eyn, section 8.2]）。

\smallskip
\noindent\textbf{提示}。回忆 Euler Beta 函数的积分表示
\begin{equation}
  \mathrm{B}(p,q)=\frac{\Gamma(p)\Gamma(q)}{\Gamma(p+q)}=\int_{0}^{1}t^{p-1}(1-t)^{q-1}\,\dd t
  \tag{3.59}
\end{equation}
以及 Euler Gamma 函数的渐近展开
\begin{equation}
  \ee^{\frac{1}{v}}\sqrt{2\pi}\,\frac{(-v)^{\frac{1}{v}+a+\frac{1}{2}}}{\Gamma(a-v^{-1})}\sim\exp\left(\sum_{m=1}^{\infty}\frac{B_{m+1}(a)}{m(m+1)v^m}\right).
  \tag{3.60}
\end{equation}
这里 $B_{m+1}(a)$ 是 Bernoulli 多项式，在 $a=0$ 与 $a=1$ 处都特殊化为 Bernoulli 数：$B_{m+1}(0)=B_{m+1}(1)=B_{m+1}$。
\end{exercise}

% [Tutor Note: (3.58) 的分歧点：$x'(z)=-f/z+1/(1-z)$，令其为零得 $z=f/(1+f)$（$f\ne0,-1$）；在该点 $x''(z)=f/z^2+1/(1-z)^2\ne0$，故为单分歧点。因此该谱曲线恰有一个分歧点，与 $\dim V=1$（表 2 第一行）相合。]

% [Tutor Note: “triple Hodge” 即三个 Hodge 类的乘积 $\Lambda(1)\Lambda(f)\Lambda(-f-1)$（$\Lambda(t)$ 在 $t=1,f,-f-1$ 处的取值之积），其中 $f$ 是 framed 拓扑顶点的 framing 参数，它出现在 (3.58) 的 $x(z)=-f\log z-\log(1-z)$ 中。原文未给出 “large framing limit” 的确切取法，只指出极限后得到计算 Hurwitz 数的 Lambert 曲线，细节见所引 [Eyn, section 8.2]。]

% [Tutor Note: 核对 (3.59)：$p=q=1$ 时 $B(1,1)=\Gamma(1)^2/\Gamma(2)=1$，右端 $\int_0^1\dd t=1$，两式一致。(3.59) 把 Euler Beta 函数写成 Euler 积分，习题 3.5 的证明用它把 $\int_0^1t^{p-1}(1-t)^{q-1}\dd t$ 型积分化为 Gamma 函数之比。]

% [Tutor Note: (3.60) 中的 $v$ 是拉丁字母 $v$，是渐近展开参数（地位相当于 (3.53)–(3.54) 中的 $u$），不是分歧点的希腊字母 $\nu$；(3.53) 中 $x(\nu)$ 的 $\nu$ 与它形状相近，抄写时务必区分。右端是 $v\to0$ 的渐近级数（一般发散），等式只在形式渐近展开的意义下成立。]

% [Tutor Note: 原文说 $B_{m+1}(a)$ “在 $a=0$ 与 $a=1$ 处都特殊化为 Bernoulli 数”：严格说这仅对 $m+1\ge2$ 成立，因为 $B_1(0)=-\frac12=B_1$ 而 $B_1(1)=+\frac12$。由于 (3.60) 的求和从 $m=1$ 开始（故 $m+1\ge2$），该断言在用到的地方无误。]

我们最后指出，拓扑递归与上同调场论之间的对应，近来在高能物理的若干领域中找到了应用，包括引力 [Ebe24; CEMR24; CEMR25; AC25] 与规范场论 [GM]。
